% =============================================================================
%  II. BACKGROUND: THE GAP BETWEEN PERFORMANCE AND CLINICAL UTILITY
% =============================================================================
\begin{frame}{II. Background: The Gap Between Predictive Performance and
              Clinical Utility}
Wang et al.~\cite{wang2025srpm} reviewed 91 real-time sepsis prediction
models and found:

\vspace{2pt}
\begin{center}\tablestyle\footnotesize
\begin{tabular}{lcc}
\toprule
\textbf{Metric} & \textbf{Internal} & \textbf{External} \\
\midrule
Median AUROC         & 0.811           & 0.783 \\
Median Utility Score & \textbf{+0.381} & \textbf{--0.164} \\
\bottomrule
\end{tabular}
\end{center}

\vspace{2pt}
\begin{itemize}\itemsep1pt\small
  \item AUROC stays almost the same, so the model looks good.
  \item But the Utility Score becomes negative, meaning the model's alerts
        are causing more problems than benefits.
  \item AUROC shows prediction ability, but Utility
        Score~\cite{reyna2020utility} shows whether the model is actually
        useful in real clinical practice.
\end{itemize}

\vspace{2pt}
\begin{alertblock}{}
  \centering The field has a failure it cannot see with its standard metrics.
\end{alertblock}
\end{frame}


% =============================================================================
%  II. BACKGROUND: LIMITATIONS OF THE CONVENTIONAL APPROACH
% =============================================================================
\begin{frame}{II. Background: Limitations of the Conventional Approach}
The traditional approach is:

\vspace{2pt}
{\centering Choose one Sepsis-3 definition $\rightarrow$ use one train/test
split $\rightarrow$ report AUROC $\rightarrow$ compare models.\par}

\vspace{4pt}
\textbf{Four key limitations:}
\vspace{-2pt}
\begin{enumerate}\itemsep1pt\footnotesize
  \item \textbf{The sepsis label can change.} Cohen et
        al.~\cite{cohen2024subtle} showed that different ways of applying
        Sepsis-3 criteria to the same database can create very different
        patient groups (867--2,178 cases). The choice of label can affect
        results as much as the choice of model.
  \item \textbf{The data split can give unrealistic results.} Guo et
        al.~\cite{guo2022temporal} found that sepsis prediction performance
        drops significantly when tested on data from a different time period.
        Random train/test splits may make models look better than they
        actually are.
  \item \textbf{AUROC does not show the full picture.} Wang et
        al.~\cite{wang2025srpm} showed that a model can maintain a good AUROC
        while its clinical usefulness decreases. A high AUROC does not always
        mean the model helps doctors.
  \item \textbf{The label may represent doctor decisions, not only patient
        deterioration.} Kamran et al.~\cite{kamran2024evaluation} suggested
        that models trained on clinical data may learn patterns of medical
        actions instead of learning the actual worsening of a patient's
        condition.
\end{enumerate}
\end{frame}


% =============================================================================
%  II. BACKGROUND: WHERE PERFORMANCE VARIATION COMES FROM
% =============================================================================
\begin{frame}{II. Background: Where Performance Variation Comes From}
{\small The traditional approach assumes that the model is the main factor
affecting performance. Other things, such as the sepsis label, data split,
and evaluation method, are treated as minor technical choices.\par}

\vspace{4pt}
\textbf{However, research shows that these choices can strongly affect the
results:}
\begin{itemize}\itemsep1pt\small
  \item The choice of sepsis label can change model performance as much as
        changing the model itself~\cite{cohen2024subtle}.
  \item The way data is divided into training and testing sets can decide
        whether the reported performance is realistic or
        misleading~\cite{guo2022temporal}.
  \item The choice of evaluation metric can determine whether a model failure
        is detected or hidden~\cite{wang2025srpm}.
\end{itemize}

\vspace{4pt}
\begin{alertblock}{}
  \centering Performance differences may come more from research decisions
  than from the AI model itself.
\end{alertblock}

\vspace{4pt}
{\small This study measures how much performance variation comes from the
model, the label, and the analyst's choices.\par}
\end{frame}
